%% ma: L12-B14-0050
%% lop: 12
%% chuong: 5
%% bai: 14
%% dang: 12.14.1
%% loai: TN4
%% muc: NB
%% dap_an: "C"
%% nguon: "To_n_12_PT_M_t_Ph_ng.pdf (Đề kiểm tra đánh giá kiến thức, Chuyên đề Phương trình mặt phẳng, Nguyễn Anh Khôi - Lê Minh Khôi), Câu 4"
%% kien_thuc: "Bài 14 (phương trình tổng quát của mặt phẳng, vectơ pháp tuyến)"
%% trang_thai: chinh-thuc
%% ghi_chu: "Tài liệu gốc không kèm đáp án; đáp án do Claude giải và kiểm bằng sympy."
\begin{ex}
Trong không gian $Oxyz$, cho mặt phẳng $(P)\colon z-2x+3=0$. Một vectơ pháp tuyến của $(P)$ là
\choice{$\vec u=(0;1;-2)$}{$\vec v=(1;-2;3)$}{\True $\vec n=(2;0;-1)$}{$\vec w=(1;-2;0)$}
\loigiai{$(P)\colon -2x+0y+z+3=0$ có vectơ pháp tuyến $(-2;0;1)$, cùng phương với $\vec n=(2;0;-1)$. Chọn C.}
\end{ex}
